\documentclass[12pt, a4paper]{article}

% --- Pacotes Básicos ---
\usepackage[utf8]{inputenc}
\usepackage[T1]{fontenc}
\usepackage[main=brazilian, english]{babel}
\usepackage{indentfirst}
\usepackage{geometry}
\geometry{left=3cm, top=3cm, right=2cm, bottom=2cm}

% --- Informações do Documento ---
\title{Aula Assíncrona: Encontro 7\\Iniciação à Prática de Extensão}
\author{Nome: Erickson Giesel Müller
\\ Fase: 6° Semestre Ciência da Computação
\\ G5 Saúde Digital}
\date{\today}

\begin{document}
\maketitle
\thispagestyle{empty} % primeira página não tem número.

O texto Informática em Saúde e Cidadania Digital aborda a evolução do desenvolvimento em informática em saúde, passando pelos primeiros usos de tecnologia digital na área da saúde, em 1960, o desenvolvimento do conceito de "informática médica" o qual define o uso de informática para auxiliar no desempenho médico armazenando estatísticas e dados que posteriormente veio a se tornar "informática biomédica", em que a informática passou a ser peça importante no estudo, análise e desenvolvimento da área.

Atualmente, a comunidade internacional adota o termo "informática em saúde" e é pensado tendo em vista o desenvolvimento de informática muito antes da ideia de computação digital que temos hoje em dia, Homer R. Warner foi um nome importante para a adoção global de dispositivos de informática em hospitais clínicos, em 1956 o cardiologista usou a computação para desenvolver um sistema de apoio à decisão clínica, que oferecia sugestões de diagnósticos para atendimentos cardiológicos.

Essa alteração no curso da história comprovou o teorema fundamental da Informática em Saúde, proposto por Charles Friedman, em que profissionais de saúde que utilizam um sistema de informação em saúde obtêm melhores resultados que profissionais que não utilizam. Esse teorema é muito mais perceptível hoje em dia, em que a quantidade de informações que os profissionais da saúde precisam consultar e interpretar ultrapassa as capacidades cognitivas de qualquer ser humano.

O autor explicou o conceito de Cidadania Digital, que surge para implantar as normas de cidadania, que sempre tivemos, mas agora em ambientes digitais. A Cidadania Digital anda lado a lado com o conceito de Inteligência Digital, conhecer as oito áreas da Inteligência Digital, e suas devidas competências, é de suma importância para o desenvolvimento de comportamentos que integram a Cidadania Digital. Encontrar e desenvolver as competências é uma tarefa que deve ser feita ao analisar com um olhar crítico os comportamentos que temos quando nos expomos ao meio digital.

\end{document}
